\chapter{第一题}

\section{题目}

并行化时我们需要注意并识别循环依赖，分析以下两个 C语言循环片段，判断其是否可以直接使用\texttt{\#pragma omp parallel for}进行并行化。如果不能，请说明原因：

片段 A：
\begin{lstlisting}[style=cleancpp]
for(i= 1; i< n; i++){ 
    A[i]= A[i-1]+ B[i]; 
}
\end{lstlisting}

片段 B：
\begin{lstlisting}[style=cleancpp]
int index= 0;
for(i= 0; i< n; i++){ 
    index=(i+ offset)% n; //offset为一个固定的常数 
    C[index]= A[i]* B[i]; 
}
\end{lstlisting}

\section{答案}
\subsection{片段A的并行化分析}

片段A\textbf{不能}直接使用\texttt{\#pragma omp parallel for}进行并行化。

原因：存在\textbf{循环携带依赖}（Loop-carried dependence）。第$i$次迭代计算$A[i]$时需要依赖第$i-1$次迭代计算出的$A[i-1]$值。这种数据依赖关系使得迭代之间无法独立执行，必须按顺序串行执行。

\subsection{片段B的并行化分析}

片段B\textbf{在理论上可以并行化}，但在实际使用\texttt{\#pragma omp parallel for}时可能存在问题。

原因：表达式 $index = (i + offset) \% n$ 实际上是对数组索引的一种置换（permutation），在数学上保证了不同的 $i$ 会映射到不同的 $index$，因此不会发生写冲突。

然而，编译器可能无法在编译时证明这种映射关系是无冲突的，因为这需要进行复杂的数学推理分析来保证不会发生冲突，因此从实际使用的角度出发，不一定能够安全地自动并行化。

\chapter{第二题}

\section{题目}

请阅读以下使用 OpenMP并行化的 C语言代码片段，假设 n非常大，线程数为 4个：
\begin{lstlisting}[style=cleancpp]
double result= 0.0;
//使用 OpenMP并行化 
#pragma omp parallel for reduction(+: result) 
for(int i= 0; i< n; i++){ 
    result+= HeavyCompute(i); 
}
\end{lstlisting}
其中 HeavyCompute(i)是一个模拟函数，i越大，计算所需的时间越长。

回答下列问题：
\begin{enumerate}
    \item 虽然这段代码计算结果是正确的，但在性能上可能会出现什么问题？
    \item OpenMP的 parallel for默认使用哪种调度策略？为什么这种策略在上述场景下效率低下？
    \item 请修改\#pragma指令，使用一种更合适的调度策略来平衡负载，提高程序运行速度。
\end{enumerate}

\section{答案}
\subsection{性能问题分析}

虽然代码计算结果正确，但在性能上会出现\textbf{负载不均衡}（Load Imbalance）问题。

由于HeavyCompute(i)的计算时间随$i$增大而增加，而默认的静态调度策略会将连续的迭代块分配给各个线程，导致处理较大$i$值的线程需要更长时间完成，而处理较小$i$值的线程会提前完成并处于空闲状态，造成CPU资源浪费。

\subsection{默认调度策略及效率问题}

OpenMP的\texttt{parallel for}默认使用\textbf{静态调度}（static scheduling）策略，通常将迭代空间均匀地划分为连续的块（chunk size约为 $\frac{n}{\text{线程数}}$）分配给各个线程。
在这种场景下效率低下的原因：
\begin{itemize}
    \item 静态调度无法适应不同迭代计算量的差异
    \item 计算量大的迭代集中在某些线程，造成这些线程成为性能瓶颈
    \item 其他线程完成工作后只能等待，无法帮助处理剩余的重负载
    \item 整体执行时间由最慢的线程决定，无法充分利用所有CPU核心
\end{itemize}

\subsection{改进的调度策略}

\textbf{方案一：使用动态调度}

修改\texttt{\#pragma}指令：

\begin{lstlisting}[style=cleancpp]
#pragma omp parallel for reduction(+: result) schedule(dynamic, chunk_size)
\end{lstlisting}

其中\texttt{chunk\_size}可以根据实际情况调整，通常设置为较小的值（如1、10或100）以获得更好的负载均衡。
动态调度会在运行时将迭代块动态分配给空闲线程，每个线程完成当前块后会请求新的任务块。\textbf{较小的chunk\_size能实现更精细的负载均衡}，但会增加调度开销；
较大的chunk\_size减少调度开销但可能降低均衡效果。建议根据HeavyCompute(i)的计算时间差异程度选择合适的chunk\_size。

\textbf{方案二：使用指导性调度}

\begin{lstlisting}[style=cleancpp]
#pragma omp parallel for reduction(+: result) schedule(guided)
\end{lstlisting}

指导性调度是一种自适应的动态调度策略，初始时分配较大的任务块以减少调度开销，随着迭代进行逐渐减小块大小以提高负载均衡。
它自动平衡调度开销和负载均衡，特别适合计算时间差异大且不规律的场景，无需手动调整chunk\_size参数。

这两种策略都能让线程在完成当前任务后动态获取新的任务，从而更好地平衡负载。

\chapter{第三题}

\section{题目}

请阅读以下代码片段并回答问题：
\begin{lstlisting}[style=cleancpp]
#include<stdio.h>
#include<omp.h> 
int main(){ 
    int i, s= 0, a= 5;
    #pragma omp parallel num_threads(3) private(i) shared(a) reduction(+:s) { 
        int tid= omp_get_thread_num(); 
        a+= tid;
        #pragma omp for 
        for(i= 1; i<= 6; i++){ 
            s+= i;
        }
    } 
    printf("s=%d\n", s); 
    printf("a=%d\n", a); 
    return 0;
}
\end{lstlisting}

回答下列问题：
\begin{enumerate}
    \item 变量 i, s, a分别是什么变量？
    \item 指出这段代码所存在的问题。
    \item 分析最终 a和 s的输出。
\end{enumerate}

\section{答案}
\subsection{变量类型分析}

\begin{itemize}
    \item \textbf{i}: 私有变量（private）。在并行区域内每个线程都有自己的副本
    \item \textbf{s}: 归约变量（reduction variable）。每个线程有私有副本，最后通过加法归约合并
    \item \textbf{a}: 共享变量（shared）。所有线程共享同一个内存位置
\end{itemize}

\subsection{代码存在的问题}

代码存在\textbf{数据竞争}（Data Race）问题：

变量$a$被声明为共享变量，但多个线程同时执行\texttt{a += tid}操作。由于这个操作不是原子的，会导致数据竞争，产生不确定的结果。

具体来说，三个线程（tid = 0, 1, 2）都会同时读取、修改和写入共享变量$a$，可能导致某些增量操作丢失。

\subsection{最终输出分析}

\textbf{s的输出}：$s = 21$

原因：归约操作保证了正确性。循环\texttt{for(i = 1; i <= 6; i++)}的总和为$1+2+3+4+5+6 = 21$，归约操作会正确地将所有线程的局部和相加。

\textbf{a的输出}：\textbf{不确定}

理论上，如果没有数据竞争，$a$应该等于$5 + 0 + 1 + 2 = 8$。但由于存在数据竞争，实际结果通常为6、7或8，取决于线程执行的交错顺序。

\chapter{第四题}

\section{题目}

某程序使用 OpenMP并行化如下循环：
\begin{lstlisting}[style=cleancpp]
#pragma omp parallel for num_threads(4) 
for(int i= 0; i< 16; i++){ 
    work(i); 
}
\end{lstlisting}
设系统采用默认调度方式，并假设每次迭代 work(i)的执行时间基本相同。

回答下列问题：
\begin{enumerate}
    \item 若 16次迭代被连续地、尽可能均匀地分配给 4个线程，请写出每个线程可能负责的迭代编号范围。
    \item 这种分配方式有什么优点？
    \item 如果不同迭代的计算量差异非常大，那么默认这种分配方式可能会带来什么问题？
\end{enumerate}

\section{答案}
\subsection{迭代分配范围}

在默认静态调度下，16次迭代被均匀分配给4个线程：

\begin{itemize}
    \item 线程0：迭代 0, 1, 2, 3（范围：0-3）
    \item 线程1：迭代 4, 5, 6, 7（范围：4-7）
    \item 线程2：迭代 8, 9, 10, 11（范围：8-11）
    \item 线程3：迭代 12, 13, 14, 15（范围：12-15）
\end{itemize}

\subsection{分配方式的优点}

这种静态连续分配的优点包括：
\begin{itemize}
    \item \textbf{良好的空间局部性}：连续的内存访问模式有利于缓存性能
    \item \textbf{低调度开销}：在进入并行区域时就确定了分配方案，无需运行时调度
    \item \textbf{可预测性}：每次运行的分配结果一致，便于调试和性能分析
    \item \textbf{适合均匀负载}：当各迭代计算量相近时，能实现良好的负载均衡
\end{itemize}

\subsection{计算量差异大时的问题}

当不同迭代的计算量差异很大时，静态分配会导致\textbf{严重的负载不均衡}：

\begin{itemize}
    \item 某些线程可能分配到计算量很大的迭代，需要很长时间完成
    \item 其他线程完成自己的轻量级迭代后只能空闲等待
    \item 整体执行时间由最慢的线程决定，无法充分利用所有CPU核心
    \item 并行效率显著降低，甚至可能比串行执行还要慢（考虑并行开销）
\end{itemize}

\chapter{第五题}

\section{题目}

阅读下面的 OpenMP程序片段：
\begin{lstlisting}[style=cleancpp]
double x= 0.0; 
double y= 0.0;
#pragma omp parallel num_threads(4) { 
    for(int i= 0; i< 1000; i++){ 
        #pragma omp atomic 
        x+= 1.0;
        #pragma omp atomic 
        y+= 1.0; 
    } 
}
\end{lstlisting}

回答下列问题：
\begin{enumerate}
    \item 这两条 atomic指令分别保护的是什么操作？
    \item 对变量 x的更新和对变量 y的更新，是否会被视为同一个"临界区"？为什么？
    \item 如果把两条 atomic改成一个 critical代码块，将两次加法都放进去，程序的并发性会发生什么变化？
\end{enumerate}

\section{答案}
\subsection{atomic指令保护的操作}

两条atomic指令分别保护以下操作：
\begin{itemize}
    \item 第一条：保护对变量$x$的读-修改-写操作（$x += 1.0$）
    \item 第二条：保护对变量$y$的读-修改-写操作（$y += 1.0$）
\end{itemize}

每个atomic指令确保对应的加法操作是原子执行的，防止多个线程同时修改同一变量导致的数据竞争。

\subsection{临界区分析}

对变量$x$的更新和对变量$y$的更新\textbf{不会}被视为同一个临界区。

原因：OpenMP的atomic指令是针对特定内存位置的。每个atomic指令只保护对其指定变量的访问。由于$x$和$y$是不同的内存位置，它们的atomic操作是独立的，可以并发执行。一个线程在执行$x += 1.0$的同时，另一个线程可以执行$y += 1.0$，两者互不影响。

\subsection{使用critical代码块的影响}

如果将两条atomic指令改为一个critical代码块：

\begin{lstlisting}[style=cleancpp]
#pragma omp critical
{
    x += 1.0;
    y += 1.0;
}
\end{lstlisting}

程序的并发性会\textbf{显著降低}：

\begin{itemize}
    \item 原来的atomic版本允许对$x$和$y$的更新并发进行
    \item critical版本强制所有线程串行执行整个代码块
    \item 即使不同线程操作不同的变量，也必须排队等待
    \item 并发度从可以同时更新两个变量降低到只能串行执行
    \item 性能会明显下降，特别是在高并发场景下
\end{itemize}

因此，在保护不同变量的情况下，使用单独的atomic指令比使用critical代码块更高效。